% P15-U1 continuation: measurable probability foundations for N-x-008.
\chapter{Measurable Probability and Random Variables}
\label{sec:p15-measurable-probability}

An infinite probability model must assign probabilities consistently to
countable combinations of events. Its numerical observables must also have
measurable inverse images, so that a question about an observed value names
an event to which the model assigns a probability. These two requirements
connect probability with measure theory.

This chapter develops the measure, distribution and random variable
foundations in Rick Durrett's \emph{Probability: Theory and Examples},
author version 5 dated 11 January 2019
\cite[secs.~1.1--1.3; app.~A.1]{DurrettPTEAuthor2019}.
The arguments assume elementary set theory, finite sums, and the order
completeness and elementary topology of the real numbers. Countable sums
of nonnegative terms are defined below. Every stated theorem has its proof
in the annex. The countable models are constructed here; the construction
of Lebesgue measure and the theory of integration have a separate boundary
at the end of the chapter.

\section{Event collections and measures}
\label{sec:p15-event-measures}

\begin{hbdefinition}{def:p15-sigma-algebra}{Measurable spaces and generated event collections}
Let \(\Omega\) be a nonempty set. Its elements are \emph{outcomes};
\(2^\Omega\) denotes its set of all subsets, and
\(A^c\defeq\Omega\setminus A\). A \emph{sigma-algebra}, also called
a \emph{sigma-field}, on \(\Omega\) is a collection
\(\mathcal F\subseteq2^\Omega\) containing \(\Omega\), closed under
complements and countable unions. Here a countable family may be finite
or countably infinite. An empty union is \(\varnothing\), and an empty
intersection of subsets of \(\Omega\) is \(\Omega\).
The pair \((\Omega,\mathcal F)\) is a \emph{measurable space}, and
members of \(\mathcal F\) are its \emph{measurable sets} or
\emph{events}. A collection containing \(\Omega\) and closed under
complements and finite unions is an \emph{algebra} of sets.

For any collection \(\mathcal C\subseteq2^\Omega\), define
\[
 \sigma(\mathcal C)\defeq
 \bigcap\{\mathcal H:\mathcal H\text{ is a sigma-algebra on }\Omega,
                              \ \mathcal C\subseteq\mathcal H\}.
\]
The family being intersected is nonempty because it contains \(2^\Omega\).
\index{sigma-algebra}\index{measurable space}\index{event!measurable}
\end{hbdefinition}

\begin{hbtheorem}{thm:p15-generated-events}{proved=annex:p15.generated-events}{Closure and generation}
Every sigma-algebra contains \(\varnothing\) and is closed under
countable intersections and differences of two of its sets. An intersection
of any nonempty family of sigma-algebras on the same set is a sigma-algebra.
Consequently \(\sigma(\mathcal C)\) is the unique smallest sigma-algebra
containing \(\mathcal C\): it contains \(\mathcal C\) and is contained
in every sigma-algebra containing \(\mathcal C\).
\end{hbtheorem}

\begin{hbexample}{def:p15-finite-cofinite-events}{Finite closure does not supply countable closure}
Let \(\mathbb N_+\defeq\{1,2,3,\ldots\}\), and let \(\mathcal A\)
contain exactly the finite and cofinite subsets of \(\mathbb N_+\),
where cofinite means having finite complement. Complements interchange
the two kinds of set. A union of two finite sets is finite; a union
containing a cofinite set is cofinite. Thus \(\mathcal A\) is an algebra.
Every singleton belongs to it, but the union of the even singletons is
neither finite nor cofinite. Hence \(\mathcal A\) is not a sigma-algebra.
Its generated sigma-algebra is \(2^{\mathbb N_+}\), because each subset
of \(\mathbb N_+\) is a countable union of singletons.
\end{hbexample}

\begin{hbdefinition}{def:p15-measure}{Nonnegative sums, measures and probability}
Write \([0,\infty]\) for the nonnegative extended real numbers, ordered
by putting \(\infty\) above every real number. Addition satisfies
\(a+\infty=\infty+a=\infty\), including \(a=\infty\).
For nonnegative terms \(a_i\in[0,\infty]\) indexed by a countable set
\(I\), set
\[
 \sum_{i\in I}a_i\defeq
 \sup\left\{\sum_{i\in J}a_i:J\subseteq I\text{ finite}\right\}.
\]
An empty finite sum is zero. For a sequence this is the supremum, or
increasing limit, of its initial partial sums: every finite set of indices
is included in an initial segment. No subtraction of infinities is defined.

A \emph{measure} on \((\Omega,\mathcal F)\) is a function
\(\mu:\mathcal F\to[0,\infty]\) such that
\(\mu(\varnothing)=0\) and, for every sequence of pairwise disjoint
events \((A_n)_{n\geq1}\),
\[
 \mu\left(\bigcup_{n\geq1}A_n\right)=\sum_{n\geq1}\mu(A_n).
\]
The last property is \emph{countable additivity}. The measure is
\emph{finite} when \(\mu(\Omega)<\infty\), and it is a
\emph{probability measure} when \(\mu(\Omega)=1\). A probability
space is a triple \((\Omega,\mathcal F,\mathbb P)\) with such a
probability measure. A measurable set of measure zero is \emph{null}.
A property with a specified event holds \emph{almost surely} if that
event has probability one.
\index{measure}\index{probability!countable additivity}
\index{almost surely}\index{null event}
\end{hbdefinition}

\begin{hbtheorem}{thm:p15-measure-laws}{proved=annex:p15.measure-laws}{Additivity, bounds and continuity}
Let \(\mu\) be a measure on \((\Omega,\mathcal F)\). All sets below
belong to \(\mathcal F\).
\begin{enumerate}
\item The measure is finitely additive on pairwise disjoint events,
      including an empty family. If \(A\subseteq B\), then
      \(\mu(B)=\mu(A)+\mu(B\setminus A)\), and
      \(\mu(A)\leq\mu(B)\). If \(\mu(B)<\infty\), then
      \(\mu(B\setminus A)=\mu(B)-\mu(A)\).
\item For any sequence \((A_n)\),
      \(\mu(\bigcup_n A_n)\leq\sum_n\mu(A_n)\).
      A countable union of null events is null.
\item If \(A_n\subseteq A_{n+1}\) and \(A=\bigcup_n A_n\), then
      \(\mu(A_n)\) increases to \(\mu(A)\), possibly to infinity.
\item If \(A_{n+1}\subseteq A_n\), \(A=\bigcap_n A_n\), and
      \(\mu(A_1)<\infty\), then \(\mu(A_n)\) decreases to
      \(\mu(A)\).
\end{enumerate}
For a probability measure, \(0\leq\mathbb P(A)\leq1\),
\(\mathbb P(A^c)=1-\mathbb P(A)\), and a countable intersection of
events of probability one has probability one. Both continuity conclusions
therefore apply to probability measures without an extra finiteness premise.
\end{hbtheorem}

\section{Constructing countable models}
\label{sec:p15-countable-models}

\begin{hbtheorem}{thm:p15-countable-atomic}{proved=annex:p15.countable-atomic}{Grouping nonnegative sums and assigning masses}
Let \(I\) be countable, \(a_i\in[0,\infty]\), and let
\((I_n)_{n\geq1}\) be pairwise disjoint subsets whose union is \(I\).
Empty members are allowed. Then
\[
 \sum_{i\in I}a_i=\sum_{n\geq1}\left(\sum_{i\in I_n}a_i\right).
\]
In particular, if \(\Omega\) is nonempty and countable and
\(p:\Omega\to[0,\infty)\), the assignment
\[
 \mu_p(A)\defeq\sum_{\omega\in A}p(\omega)
       \qquad(A\subseteq\Omega)
\]
is a measure on \((\Omega,2^\Omega)\). If
\(\sum_{\omega\in\Omega}p(\omega)=1\), it is a probability measure
and \(p\) is its \emph{probability mass function}. Conversely, every
probability measure on \((\Omega,2^\Omega)\) arises uniquely in this
way, with \(p(\omega)=\mathbb P(\{\omega\})\).
\end{hbtheorem}

The first conclusion makes rearrangement and grouping legitimate for
nonnegative series, including divergent ones. It asserts no rearrangement
rule for a signed series. A probability space obtained by this mass
assignment will be called a \emph{countable atomic model}. The atomic
construction recovers the finite
models of Chapter~\ref{sec:p15-finite-probability} with the same weights
and event probabilities.

\begin{hbexample}{def:p15-countable-waiting}{A countably infinite waiting model}
Take \(\Omega=\mathbb N_+\), with all subsets measurable, and put
\(p(n)=2^{-n}\). The finite geometric identity
\(\sum_{n=1}^m2^{-n}=1-2^{-m}\), obtained by subtracting half the sum
from the sum, gives total mass one as \(m\to\infty\).
Interpret \(n\) as the recorded trial on which a response first occurs.
The assigned masses define this model directly; an infinite sequence of
independent trials has not been constructed as an underlying experiment.

The event \(A_m=\{n:n\geq m\}\) says the response has not occurred
before trial \(m\). Its probability is
\[
 \mathbb P(A_m)=2^{1-m},\qquad
 A_{m+1}\subseteq A_m,\qquad \bigcap_{m\geq1}A_m=\varnothing.
\]
Indeed, a fixed \(n\) fails membership once \(m=n+1\). Thus events
can all be nonempty while their intersection is empty. Their probabilities
decrease to zero, as continuity from above requires. The probabilities
of \(\{1,\ldots,m\}\) increase to one by continuity from below.

The even outcomes have probability
\(\sum_{k\geq1}2^{-2k}=(1/4)/(1-1/4)=1/3\); the odd outcomes
therefore have probability \(2/3\). Both evaluations use nonnegative
sums, with no limiting independence premise.
\end{hbexample}

\begin{hbwarning}{def:p15-infinite-measure-continuity}{The finite measure premise cannot be dropped}
On the same countable event space, the weights \(p(n)=1\) construct
\emph{counting measure}: a finite set has measure equal to its number of
elements, and an infinite set has measure \(\infty\). Every tail
\(A_m=\{m,m+1,\ldots\}\) has measure \(\infty\), while its
decreasing intersection is empty and has measure zero. Consequently
\(\lim_m\mu(A_m)=\infty\ne0=\mu(\bigcap_m A_m)\).
The failed subtraction would have involved \(\infty-\infty\).

There is also no uniform probability mass on \(\mathbb N_+\).
If all singletons had the same mass \(c>0\), sufficiently many of them
would have total mass greater than one. If \(c=0\), countable additivity
would give total mass zero. On the finite-cofinite algebra of
Example~\ref{def:p15-finite-cofinite-events}, assigning zero to finite
sets and one to cofinite sets is finitely additive: two disjoint sets
cannot both be cofinite, and adding a finite set to a disjoint cofinite
set leaves a cofinite set. This assignment has no probability measure
extension to \(2^{\mathbb N_+}\), since every singleton would remain null.
\end{hbwarning}

\section{Measurable observations and their limits}
\label{sec:p15-measurable-observations}

\begin{hbdefinition}{def:p15-borel-sets}{Borel sets}
For an integer \(d\geq1\), give \(\mathbb R^d\) its usual Euclidean
topology. A set is open if each of its points has a sufficiently small
open Euclidean ball contained in the set. The \emph{Borel sigma-algebra}
\(\mathcal B(\mathbb R^d)\) is generated by the open sets. Its members
are \emph{Borel sets}. In one dimension use \(\mathcal B(\mathbb R)\).
Write \(\mathbb Q\) for the rational numbers.
\index{Borel set}
\end{hbdefinition}

\begin{hbtheorem}{thm:p15-borel-generators}{proved=annex:p15.borel-generators}{Countable generators of Borel sets}
The Borel sigma-algebra on \(\mathbb R^d\) is generated by the open
rectangles \(\prod_{j=1}^d(a_j,b_j)\) whose endpoints are rational
and satisfy \(a_j<b_j\). In dimension one it is also generated by
\(\{(-\infty,q]:q\in\mathbb Q\}\), and by
\(\{(-\infty,a]:a\in\mathbb R\}\).
\end{hbtheorem}

\begin{hbdefinition}{def:p15-measurable-map}{Random elements and inverse-image events}
Let \((\Omega,\mathcal F)\) and \((S,\mathcal S)\) be measurable
spaces. For a function \(X:\Omega\to S\), its \emph{inverse image}
of \(B\subseteq S\) is
\[
 X^{-1}(B)\defeq\{\omega\in\Omega:X(\omega)\in B\}.
\]
This notation does not assert that \(X\) is invertible. The map is
\emph{measurable} when \(X^{-1}(B)\in\mathcal F\) for every
\(B\in\mathcal S\). On a probability space it is called an
\(S\)-valued \emph{random element}. When
\((S,\mathcal S)=(\mathbb R,\mathcal B(\mathbb R))\), it is a
\emph{real random variable}; for \(\mathbb R^d\), it is a
\emph{random vector}. All functions here are defined on all of \(\Omega\).
Write \(\{X\in B\}\) for \(X^{-1}(B)\) and
\(\{X\leq a\}\) for the corresponding inverse image of a ray.
For a function \(f\) on \(S\), write \(f(X)\defeq f\circ X\).
For the fixed target sigma-algebra \(\mathcal S\), set
\[
 \sigma(X)\defeq\{X^{-1}(B):B\in\mathcal S\}.
\]
\index{measurable map}\index{random variable!measurable}
\index{sigma-algebra!generated by a random element}
\end{hbdefinition}

\begin{hbtheorem}{thm:p15-measurability-criterion}{proved=annex:p15.measurability-criterion}{Testing measurability and recording information}
If \(\mathcal S=\sigma(\mathcal C)\), then \(X:\Omega\to S\)
is measurable exactly when \(X^{-1}(C)\in\mathcal F\) for each
\(C\in\mathcal C\). In particular a real function is a random
variable exactly when \(\{X\leq q\}\in\mathcal F\) for every
rational \(q\).

The collection \(\sigma(X)\) is a sigma-algebra and is the smallest
one on \(\Omega\) that makes \(X\) measurable. It is generated by
\(\{X^{-1}(C):C\in\mathcal C\}\) whenever
\(\mathcal S=\sigma(\mathcal C)\). If \(X\) is measurable for
\(\mathcal F\), then \(\sigma(X)\subseteq\mathcal F\).
\end{hbtheorem}

Every function from the outcome space of a countable atomic model is
measurable for any specified target sigma-algebra: each inverse image is
a subset of that outcome space, and all its subsets are events.

\begin{hbexample}{def:p15-coarse-observation}{A coarse observation has fewer measurable distinctions}
In the waiting model, let \(Y(n)=1\) on the even outcomes and \(Y(n)=0\)
on the odd outcomes. Give \(\{0,1\}\) all its subsets as measurable
sets. Then \(\sigma(Y)\) consists of exactly the empty set, the whole
space, the even set and the odd set. The function \(X(n)=n\) is
measurable for \(2^{\mathbb N_+}\), but is not measurable for
\(\sigma(Y)\): the inverse image \(\{X\leq1\}=\{1\}\) is
absent from that four-set collection. Knowing the parity does not record
the full waiting time. A sigma-algebra contained in \(\mathcal F\) is
called a \emph{sub-sigma-algebra}; it will specify the information retained
in general conditional expectation.
\end{hbexample}

\begin{hbdefinition}{def:p15-simple-functions}{Indicators and simple functions}
For \(A\in\mathcal F\), its \emph{indicator} \(\mathbf1_A\) is one
on \(A\) and zero on \(A^c\). A real function is \emph{simple} when
it is measurable and has finite range. Equivalently, as proved below,
it has a representation \(\sum_{j=1}^m c_j\mathbf1_{A_j}\), with
real constants \(c_j\), measurable sets \(A_j\), and \(m\geq0\).
The sets need not be disjoint in this representation. The empty sum is
the zero function. Operations on real functions are pointwise.
\index{simple function}\index{indicator of an event}
\end{hbdefinition}

\begin{hbtheorem}{thm:p15-measurable-operations}{proved=annex:p15.measurable-operations}{Composition, vectors and finite operations}
The composite of measurable maps is measurable. A map into
\(\mathbb R^d\) is measurable exactly when all its coordinate
functions are real random variables. Every continuous map between
Euclidean spaces is Borel measurable. Therefore continuous functions
of finitely many real random variables are real random variables;
in particular, sums, products, absolute values, maxima and minima are
measurable. Indicators are measurable, and the two descriptions of
simple functions in Definition~\ref{def:p15-simple-functions} agree.
\end{hbtheorem}

\begin{hbdefinition}{def:p15-extended-random-variables}{Extended values and pointwise limits}
Let \(\overline{\mathbb R}\defeq[-\infty,\infty]\), ordered by
adjoining a least and a greatest element to \(\mathbb R\). Equip it
with the sigma-algebra generated by \([-\infty,a]\), for real \(a\).
An \emph{extended real random variable} is a measurable map into this
space; by the generator criterion this means that all events
\(\{Z\leq a\}\), \(a\in\mathbb R\), are measurable.
Suprema and infima in this ordered set always exist, with
\(\sup\varnothing=-\infty\) and \(\inf\varnothing=\infty\).
For a sequence \((Z_n)\), define pointwise
\[
 \liminf_n Z_n\defeq\sup_{k\geq1}\inf_{n\geq k}Z_n,
 \qquad
 \limsup_n Z_n\defeq\inf_{k\geq1}\sup_{n\geq k}Z_n.
\]
A sequence converges almost surely to a specified real random variable
\(X\) when the event on which its real values converge to \(X\)
has probability one. Measurability of that event is supplied next.
\index{random variable!extended real}\index{convergence!almost surely}
\end{hbdefinition}

\begin{hbtheorem}{thm:p15-measurable-limits}{proved=annex:p15.measurable-limits}{Limits and approximation by simple functions}
For extended real random variables \((Z_n)_{n\geq1}\), the functions
\(\inf_n Z_n\), \(\sup_n Z_n\), \(\liminf_n Z_n\) and
\(\limsup_n Z_n\) are extended real random variables. If real random
variables converge at every outcome to a finite value, their limit is
a real random variable. For real random variables \((X_n)\) and \(X\),
\[
 \{X_n\to X\}=
 \bigcap_{k\geq1}\bigcup_{N\geq1}\bigcap_{n\geq N}
       \{|X_n-X|<1/k\}
\]
is a measurable event.

Every nonnegative extended real random variable \(Z\) admits a
sequence of nonnegative simple functions increasing pointwise to \(Z\).
One explicit choice, for \(n\geq1\), is
\[
 s_n\defeq
 \sum_{k=0}^{n2^n-1}\frac{k}{2^n}
    \mathbf1_{\{k2^{-n}\leq Z<(k+1)2^{-n}\}}
       +n\mathbf1_{\{Z\geq n\}}.
\]
This formula defines a finite real function even where \(Z=\infty\).
\end{hbtheorem}

For example, the constant functions \(Z_n(\omega)=n\) have extended
limit \(\infty\). Their limit is measurable but is not real valued.
The limit theorem does not justify subtracting the limiting upper and
lower values: both may be infinite. The displayed event formula uses
only finite differences of real random variables.

\section{Laws and cumulative distributions}
\label{sec:p15-laws-distributions}

\begin{hbdefinition}{def:p15-pushforward-law}{The law of a random element}
On \((\Omega,\mathcal F,\mathbb P)\), let
\(X:\Omega\to(S,\mathcal S)\) be measurable. Its \emph{law}, or
\emph{distribution}, is the set function
\[
 \mathbb P_X(B)\defeq\mathbb P(X^{-1}(B))
                  \qquad(B\in\mathcal S).
\]
This is the \emph{pushforward} of \(\mathbb P\) by \(X\).
The expression \(\mathbb P(X\in B)\) abbreviates the same number.
If \(\{x\}\in\mathcal S\), its probability
\(p_X(x)\defeq\mathbb P_X(\{x\})=\mathbb P(X=x)\) is the
\emph{point mass} at \(x\). The law is a measure on measurable value
sets; a point mass is a number. Two random elements with the same
measurable value space have the \emph{same law} when these measures agree
on every measurable set, even if their underlying probability spaces differ.
\index{probability!law}\index{pushforward measure}\index{point mass}
\end{hbdefinition}

\begin{hbtheorem}{thm:p15-pushforward-laws}{proved=annex:p15.pushforward-laws}{Transporting probabilities through observations}
The law \(\mathbb P_X\) is a probability measure on \((S,\mathcal S)\).
If \(f:(S,\mathcal S)\to(T,\mathcal T)\) is measurable, then
\[
 \mathbb P_{f(X)}(D)=\mathbb P_X(f^{-1}(D))\qquad(D\in\mathcal T).
\]
If the source space is a countable atomic model with weights \(p\),
then \(\mathbb P_X(B)=\sum_{\omega:X(\omega)\in B}p(\omega)\).
If \(S\) is countable and \(\mathcal S=2^S\), then
\[
 \mathbb P_X(B)=\sum_{x\in B}p_X(x),\qquad
 \sum_{x\in S}p_X(x)=1.
\]
This last conclusion holds whether or not the source space is countable.
\end{hbtheorem}

In the waiting model, the law of \(X(n)=n\) has mass \(2^{-n}\) at
the positive integer \(n\). For the parity observation, grouping its
fibers gives \(\mathbb P_Y(\{1\})=1/3\) and
\(\mathbb P_Y(\{0\})=2/3\). For the real observation \(V(n)=1/n\),
the point \(1/n\) has mass \(2^{-n}\), while zero has mass zero even
though these possible values approach zero. A limit point of a value set
need not be an attained value or have positive point mass.

\begin{hbdefinition}{def:p15-pi-lambda}{Collections used to determine a measure}
A \emph{pi-system} on a set \(S\) is a collection of subsets closed
under intersections of two members. A \emph{lambda-system} on \(S\)
contains \(S\), is closed under differences \(B\setminus A\) when
\(A\subseteq B\) are members, and is closed under increasing countable
unions. An increasing union means \(\bigcup_n B_n\) with
\(B_n\subseteq B_{n+1}\).
\index{pi-system}\index{lambda-system}
\end{hbdefinition}

\begin{hbtheorem}{thm:p15-measure-uniqueness}{proved=annex:p15.measure-uniqueness}{A generating pi-system determines a probability measure}
If a lambda-system contains a pi-system \(\mathcal C\), then it
contains \(\sigma(\mathcal C)\). Consequently two probability
measures on \((S,\sigma(\mathcal C))\) that agree on every member
of the pi-system \(\mathcal C\) agree on all measurable sets.
\end{hbtheorem}

\begin{hbdefinition}{def:p15-cumulative-distribution}{Cumulative distribution function}
For a real random variable \(X\), its \emph{cumulative distribution
function} is
\(F_X(a)\defeq\mathbb P(X\leq a)\), for \(a\in\mathbb R\).
For a nondecreasing bounded real function \(F\), write
\(F(a-)\defeq\sup_{t<a}F(t)\); this is its left limit.
Right continuity at \(a\) means \(\lim_{t\downarrow a}F(t)=F(a)\).
\index{cumulative distribution function}
\end{hbdefinition}

\begin{hbtheorem}{thm:p15-cdf-laws}{proved=annex:p15.cdf-laws}{Limits, intervals and jumps of a distribution}
The function \(F_X\) is nondecreasing, takes values in \([0,1]\),
is right continuous, and satisfies
\[
 \lim_{a\to-\infty}F_X(a)=0,\qquad
 \lim_{a\to\infty}F_X(a)=1.
\]
For real \(a<b\),
\[
\begin{aligned}
 \mathbb P(a<X\leq b)&=F_X(b)-F_X(a),\\
 \mathbb P(X<a)&=F_X(a-),\\
 p_X(a)&=F_X(a)-F_X(a-).
\end{aligned}
\]
The discontinuities of \(F_X\) are exactly the points with positive
point mass, and this set is at most countable. Two real random variables
have the same law exactly when their cumulative distribution functions agree.
\end{hbtheorem}

\begin{hbexample}{def:p15-waiting-cdf}{Reading the jumps in a countable law}
For the waiting variable, \(F_X(a)=0\) when \(a<1\). For the unique
integer \(m\geq1\) such that \(m\leq a<m+1\),
\(F_X(a)=1-2^{-m}\). At \(a=m\), its left limit is
\(1-2^{-(m-1)}\), with value zero for the left limit at \(m=1\).
Thus its jump is \(2^{-m}=\mathbb P(X=m)\). The cumulative function
is right continuous, and its jump at an outcome is its point mass.
The distribution is \emph{discrete}: it gives probability one to a
countable Borel set. The same term applies to any real law with that property.
\end{hbexample}

\begin{proofboundary}{def:p15-measure-integration-boundary}{The next measure and integration prerequisites}
The measurable space \((\mathbb R,\mathcal B(\mathbb R))\) specifies
allowable events; it does not by itself specify a probability measure.
This chapter has constructed all-subset countable models and the laws of
given measurable observations. The converse assertion that every
nondecreasing, right-continuous function with the displayed endpoint limits
is a cumulative distribution requires a measure existence argument.
Durrett's construction uses Lebesgue measure on \((0,1)\)
\cite[thm.~1.2.2, p.~11]{DurrettPTEAuthor2019}.
That measure must first be constructed from interval lengths through a
measure extension theorem, with its existence proof and hypotheses
\cite[thms.~1.1.4, 1.1.9; app.~A.1]{DurrettPTEAuthor2019}.
The uniqueness theorem above supplies no such existence claim.

\Cref{sec:p15-integration-expectation} constructs simple and nonnegative
integrals, signed integrability, Fatou's lemma, dominated convergence and
expectation with their precise hypotheses
\cite[secs.~1.4--1.6]{DurrettPTEAuthor2019}.
\Cref{sec:p15-radon-nikodym} proves when a law has a density relative
to a specified reference measure and supplies a continuous singular
distribution. A density and the point mass \(P(X=a)\) have their
different declared meanings.
\Cref{sec:p15-conditional-expectation,sec:p15-regular-conditional-laws}
construct scalar conditional expectations and regular conditional laws,
using the measurable structures developed here.
\end{proofboundary}
